% preamble.tex --- shared by all logbook documentation files
% Engine: pdflatex

\documentclass[11pt, a4paper]{report}

% --- Encoding and fonts ---
\usepackage[T1]{fontenc}
\usepackage[utf8]{inputenc}
\usepackage{lmodern}

% --- Page geometry ---
\usepackage[
    top=2.5cm, bottom=2.5cm,
    left=2.8cm, right=2.8cm
]{geometry}

% --- Colours ---
\usepackage[dvipsnames]{xcolor}
\definecolor{pionorange}{RGB}{255,152,0}
\definecolor{codeblue}{RGB}{0,80,160}
\definecolor{notebg}{RGB}{255,248,230}
\definecolor{warningbg}{RGB}{255,235,235}
\definecolor{codebg}{RGB}{245,245,245}
\definecolor{tabhead}{RGB}{230,230,230}

% --- Hyperlinks ---
\usepackage[
    colorlinks=true,
    linkcolor=codeblue,
    urlcolor=codeblue,
    citecolor=codeblue,
    pdftitle={Logbook Entry Generator Documentation},
    pdfauthor={PionCT Experiment}
]{hyperref}

% --- Tables ---
\usepackage{booktabs}
\usepackage{array}
\usepackage{tabularx}
\usepackage{longtable}

% --- Code listings ---
\usepackage{listings}
\lstset{
    basicstyle=\small\ttfamily,
    backgroundcolor=\color{codebg},
    frame=single,
    framerule=0.4pt,
    rulecolor=\color{gray!40},
    breaklines=true,
    breakatwhitespace=true,
    showstringspaces=false,
    tabsize=2,
    xleftmargin=6pt,
    xrightmargin=6pt,
    aboveskip=8pt,
    belowskip=6pt,
}
\lstdefinelanguage{JSON}{
    basicstyle=\small\ttfamily,
    morestring=[b]",
    stringstyle=\color{codeblue},
    literate=
        *{:}{{{\color{black}{:}}}}{1}
         {\{}{{{\color{black}{\{}}}}{1}
         {\}}{{{\color{black}{\}}}}}{1}
         {[}{{{\color{black}{[}}}}{1}
         {]}{{{\color{black}{]}}}}{1},
}
\lstdefinelanguage{bash}{
    basicstyle=\small\ttfamily,
    morecomment=[l]{\#},
    commentstyle=\color{gray},
}
\lstdefinelanguage{javascript}{
    basicstyle=\small\ttfamily,
    keywords={const,let,var,function,return,if,else,for,await,async,true,false},
    keywordstyle=\color{codeblue},
    morestring=[b]',
    morestring=[b]",
    stringstyle=\color{BrickRed},
    morecomment=[l]{//},
    commentstyle=\color{gray},
}

% --- Inline code ---
\newcommand{\code}[1]{\texttt{\small#1}}

% --- Note / Warning boxes ---
\usepackage{tcolorbox}
\tcbuselibrary{skins, breakable}

\newtcolorbox{notebox}{
    enhanced, breakable,
    colback=notebg,
    colframe=pionorange,
    leftrule=4pt, rightrule=0pt, toprule=0pt, bottomrule=0pt,
    arc=0pt, outer arc=0pt,
    left=8pt, right=8pt, top=4pt, bottom=4pt,
    fontupper=\small,
}
\newtcolorbox{warningbox}{
    enhanced, breakable,
    colback=warningbg,
    colframe=red!70!black,
    leftrule=4pt, rightrule=0pt, toprule=0pt, bottomrule=0pt,
    arc=0pt, outer arc=0pt,
    left=8pt, right=8pt, top=4pt, bottom=4pt,
    fontupper=\small,
}

% --- Section styling ---
\usepackage{titlesec}
\titleformat{\chapter}[block]
    {\normalfont\Large\bfseries\color{pionorange}}
    {\thechapter.}{0.8em}{}
    [\vspace{-0.3em}\color{pionorange}\rule{\linewidth}{1.5pt}]
\titlespacing{\chapter}{0pt}{0pt}{12pt}

\titleformat{\section}
    {\normalfont\large\bfseries\color{black}}
    {\thesection}{0.8em}{}
\titleformat{\subsection}
    {\normalfont\normalsize\bfseries}
    {\thesubsection}{0.8em}{}

% --- Header / footer ---
\usepackage{fancyhdr}
\pagestyle{fancy}
\fancyhf{}
\fancyhead[L]{\small\color{gray}\leftmark}
\fancyhead[R]{\small\color{gray}Logbook Entry Generator}
\fancyfoot[C]{\small\thepage}
\renewcommand{\headrulewidth}{0.4pt}
\renewcommand{\chaptermark}[1]{\markboth{#1}{}}

% --- Misc ---
\usepackage{enumitem}
\setlist[itemize]{itemsep=2pt, topsep=4pt}
\setlist[enumerate]{itemsep=2pt, topsep=4pt}
\usepackage{parskip}
\setlength{\parskip}{6pt}
\setlength{\parindent}{0pt}


\begin{document}

\title{
    \vspace{2cm}
    {\color{pionorange}\rule{\linewidth}{2pt}}\\[0.6em]
    {\Huge\bfseries Logbook Entry Generator}\\[0.3em]
    {\Large User Guide}\\[0.4em]
    {\color{pionorange}\rule{\linewidth}{2pt}}\\[1em]
    {\large PionCT Experiment --- Jefferson Lab}
    \vspace{1cm}
}
\author{}
\date{\today}
\maketitle
\thispagestyle{empty}

\tableofcontents
\newpage

% ============================================================
\chapter{Introduction}
% ============================================================

This guide is for shift workers using the logbook tool during an experiment.
No installation, no accounts, and no configuration are needed on your part ---
just open the HTML file in a browser and start logging.

% ============================================================
\chapter{Opening the Tool}
% ============================================================

Open \code{logbook\_generator.html} in any modern browser (Chrome, Firefox,
Edge). No internet connection is required to use the form itself. Internet is
only needed if you want to send run data to Google Sheets.

% ============================================================
\chapter{Automatic Saving}
% ============================================================

Every field you fill in is saved to your browser's local storage automatically
as you type. If you close the tab, refresh the page, or the browser crashes,
your data will be restored when you reopen the file \textbf{in the same browser
on the same machine}.

\begin{warningbox}
\textbf{Important:} Data is saved per-browser, per-machine. If you switch to a
different browser or computer, your saved session will not be there. Use
\textbf{Import from Previous HTML} (Chapter~\ref{ch:import}) to move a session
between machines.
\end{warningbox}

The autosave timestamp is shown at the top of the page under Session Management.

% ============================================================
\chapter{Session Management}
% ============================================================

At the top of the page you have two session controls:

\begin{itemize}
    \item \textbf{Import from Previous HTML} --- paste the HTML output from a
          previous entry to continue editing it. Useful for carrying over run
          data to the next shift.
    \item \textbf{Clear All \& Start Fresh} --- wipes everything and reloads a
          blank form. \textbf{This cannot be undone.}
\end{itemize}

% ============================================================
\chapter{Filling in the Form}
% ============================================================

\section{Summary}

Free-text description of the shift. Write whatever is relevant: what was
accomplished, any issues, beam status, etc.

\section{Quick Links}

Paste URLs for the standard shift documents (run plan, checklist, beamline
screenshots, etc.). Links left blank are simply omitted from the generated
output.

\section{Additional Links (optional)}

Click \textbf{+ Add Additional Link} to add any extra URLs with custom labels.
Click \textbf{Remove} on an entry to delete it.

\section{Shift Log Entries}

Click \textbf{+ Add Shift Entry} to add a timestamped log line.

\begin{itemize}
    \item \textbf{Time} --- enter the time manually in \code{HH:MM} (24-hour)
          format, or click \textbf{Now} to fill it from the system clock.
    \item \textbf{Logbook Entry} --- free text describing what happened at that
          time.
\end{itemize}

Click \textbf{Remove} on an entry to delete it. The numbering
(Shift Entry \#1, \#2\ldots) is assigned at creation and does not change if
you delete other entries.

\section{Run Summary}

Click \textbf{+ Add Run Entry} for each run taken during the shift.

\subsection{Fields filled by you}

\begin{center}
\begin{tabular}{@{}lp{8cm}@{}}
\toprule
\textbf{Field} & \textbf{Notes} \\
\midrule
Run Number          & The DAQ run number \\
Start Time          & \code{HH:MM} format, 24-hour, set by hand from the DAQ \\
Stop Time           & \code{HH:MM} format, 24-hour, set by hand from the DAQ \\
Target              & Free text \\
Beam Current (uA)   & \\
HMS P {[}GeV{]}     & \\
SHMS P {[}GeV{]}    & \\
HMS Angle {[}deg{]} & \\
SHMS Angle {[}deg{]}& \\
Number of Events (M)& \\
Trigger Rate {[}Hz{]}& \\
Live Time [\%]      & \\
DAQ Config          & \\
Comments            & \\
\bottomrule
\end{tabular}
\end{center}

\subsection{Fields filled automatically}

\begin{center}
\begin{tabularx}{\linewidth}{@{}lX@{}}
\toprule
\textbf{Field} & \textbf{How} \\
\midrule
Total Time (min) &
    Calculated from Stop $-$ Start as soon as both are filled.
    Handles midnight crossings correctly.
    If either time is missing or malformed, it is left unchanged. \\[4pt]
Date &
    Filled at the moment you click \textbf{Generate HTML} or
    \textbf{Send to Google Sheet} (local date, \code{YYYY-MM-DD}). \\
\bottomrule
\end{tabularx}
\end{center}

% ============================================================
\chapter{Google Sheets Export}
% ============================================================

If your experiment coordinator has set up Google Sheets integration, you will
see a \textbf{Google Sheets Export Settings} section. Paste the Web App URL and
shared secret there once --- they are remembered by the browser.

Each run entry has:

\begin{itemize}
    \item A \textbf{coloured dot} (LED indicator) next to its label:
          \begin{itemize}
              \item \textbf{Red} --- not yet sent to the sheet.
              \item \textbf{Green} --- successfully sent (persists across page
                    reloads).
          \end{itemize}
    \item A \textbf{Send to Google Sheet} button --- sends that specific run's
          data.
    \item An \textbf{Export All Runs to Google Sheet} button (in the settings
          section) --- sends every run in sequence.
\end{itemize}

The dot state is saved alongside your form data, so reopening the page will
still show which runs have been exported.

\begin{notebox}
\textbf{Note:} The dot reflects whether \emph{this tool} successfully sent the
data. It does not check whether the row actually landed in the sheet. If you
have doubts, verify directly in the spreadsheet.
\end{notebox}

% ============================================================
\chapter{Generating the Logbook Entry}
% ============================================================

\begin{enumerate}
    \item Click \textbf{Generate HTML}.
    \item The raw HTML appears in the output box and a preview renders below it.
    \item Click \textbf{Copy to Clipboard}.
    \item In the JLab logbook, open the summary entry you want to update.
    \item Set the text format to \textbf{HTML}.
    \item Select all existing text (\texttt{Ctrl+A}) and paste (\texttt{Ctrl+V}).
\end{enumerate}

% ============================================================
\chapter{Importing a Previous Entry}
\label{ch:import}
% ============================================================

To continue editing a previous shift's entry:

\begin{enumerate}
    \item Open the logbook entry in a browser and view its HTML source (or copy
          it from the logbook's HTML editor).
    \item Click \textbf{Import from Previous HTML} in the tool.
    \item Paste the HTML into the text box and click \textbf{Import}.
\end{enumerate}

The tool will restore the summary, links, shift log, and run table from that
entry. The Google Sheets sent/not-sent dot state is \textbf{not} restored on
import (the tool has no way to know what was previously sent), so all imported
runs will show red dots.

% ============================================================
\chapter{Tips}
% ============================================================

\begin{itemize}
    \item Total Time auto-fills only when both Start and Stop are valid
          \code{HH:MM} values. If you type a non-standard format
          (e.g.\ \code{930} instead of \code{09:30}) it will be left as-is.
    \item You can manually override Total Time by typing directly into that
          field. It will recalculate if you later change Start or Stop.
    \item The run entry number (Run Entry \#N) is permanent --- it does not
          renumber if you delete other entries. This is intentional so numbers
          stay stable during a shift.
\end{itemize}

\end{document}
